\documentclass[journal,onecolumn]{IEEEtran}
% Supplementary material for tnnls_operator.tex. Build the main paper first: the references to
% its sections, tables and equations are read from tnnls_operator.aux.

\usepackage[T1]{fontenc}
\usepackage{amsmath,amssymb}
\usepackage{booktabs}
\usepackage{graphicx}
\usepackage[numbers,sort&compress]{natbib}
\renewcommand{\bibfont}{\footnotesize}
\usepackage{url}
\usepackage{xr}
\externaldocument{tnnls_operator}
\usepackage[hidelinks]{hyperref}

\newcommand{\KL}{\mathrm{KL}}
\renewcommand{\thesection}{S-\Roman{section}}
\renewcommand{\thetable}{S-\Roman{table}}
\renewcommand{\theequation}{S\arabic{equation}}

\begin{document}

\title{Supplementary Material for ``Equal Weight Error, Unequal Effect: Quantization Sensitivity
at the Consuming Operator''}

\author{Andrew~H.~Bond,~\IEEEmembership{Senior~Member,~IEEE}}

\maketitle

Sections, tables and equations without the prefix S refer to the main paper.

\section{Decision-level agreement and a reference perturbation}\label{supp:agree}

Let $f$ be the base model and $g$ a quantized model, evaluated on $N$ items with labels $y_i$.
We count regressions and recoveries,
\begin{equation}
\begin{aligned}
r &= \frac{1}{N}\sum_i \mathbf{1}[f(x_i)=y_i,\ g(x_i)\neq y_i],\\
c &= \frac{1}{N}\sum_i \mathbf{1}[f(x_i)\neq y_i,\ g(x_i)=y_i],
\end{aligned}
\end{equation}
whose sum $r+c$ is the flip count of \citet{dutta2024accuracy}, and whose difference $r-c$ is
the accuracy change. Correctness Agreement \citep{rababah2026illusion} equals the base
accuracy minus $r$, so it falls with every regression and cannot see recoveries. The
prediction-level agreement $A = \frac{1}{N}\sum_i \mathbf{1}[f(x_i)=g(x_i)]$ asks whether the
two models give the same answer, right or wrong, and its complement $D = 1-A$ corresponds to
the AllFlips rate of \citet{dutta2024accuracy}.

We also measure a reference-perturbation disagreement $D_0$, the disagreement between two
8-bit requantizations of $f$ that differ only by a seeded Gaussian dither whose standard
deviation is a fraction of a quantization step, which we call its amplitude. $D_0$ measures sensitivity to that procedure and grows with the dither amplitude, as
reported below. It is not an amount of disagreement that the items carry regardless of
quantization, and we do not subtract it as one. We report $D - D_0$ only as a descriptive
comparison, with $D_0$ taken at an amplitude of 0.02 of a step.

\paragraph{Setup.} The task is next-token prediction on the WikiText-2 test split
\citep{merity2017pointer}, in sequences of 256 tokens. An item is one position, and it is
correct when the top-1 token equals the true next token. Every linear layer inside the
transformer blocks is fake-quantized with a symmetric absmax quantizer per output channel, and
the embeddings and language-model head are left in full precision. We run
Qwen2.5-1.5B-Instruct \citep{qwen2025qwen25} on 6{,}120 positions, Mistral-7B-v0.1
\citep{jiang2023mistral} on 12{,}240 and Llama-3.2-3B \citep{grattafiori2024llama3} on
8{,}160. Positions within a sequence are dependent, so these counts overstate the effective
sample size. Every interval in this section is a 95\% percentile bootstrap interval over 2{,}000 resamples
of whole sequences, 24, 48 and 32 of them for the three models. $D - D_0$ is resampled in pairs,
with the same sequences for both terms. The sequences are the first chunks of the split, not a
random sample of it.

\begin{table}[h]
\centering
\caption{Accuracy understates how many predictions a quantized model changes. At 4 bits the
fraction of changed predictions, $1-A$, is several times the accuracy drop on all three models,
and the intervals do not overlap. Brackets are 95\% intervals from resampling whole sequences.
$\Delta$acc is base accuracy minus quantized accuracy.}
\label{tab:agree}
\footnotesize
\setlength{\tabcolsep}{4pt}
\begin{tabular}{llrlrrll}
\toprule
model & bits & accuracy & $\Delta$acc & $r$ & $c$ & changed, $1-A$ & $D - D_0$ \\
\midrule
Qwen2.5-1.5B & base & 0.462 & & & & & \\
 & 8 & 0.462 & 0.000 [$-$0.002, 0.001] & 0.003 & 0.004 & 0.022 [0.018, 0.026] & 0.008 [0.004, 0.012] \\
 & 4 & 0.367 & 0.095 [0.085, 0.106] & 0.131 & 0.036 & 0.409 [0.388, 0.429] & 0.395 [0.376, 0.413] \\
 & 3 & 0.002 & 0.459 [0.439, 0.481] & 0.460 & 0.001 & 0.996 [0.995, 0.998] & 0.982 [0.979, 0.985] \\
\midrule
Mistral-7B & base & 0.537 & & & & & \\
 & 8 & 0.538 & 0.000 [$-$0.001, 0.000] & 0.002 & 0.002 & 0.012 [0.010, 0.013] & 0.005 [0.003, 0.007] \\
 & 4 & 0.495 & 0.043 [0.037, 0.048] & 0.072 & 0.029 & 0.232 [0.224, 0.240] & 0.226 [0.217, 0.234] \\
 & 3 & 0.000 & 0.537 [0.519, 0.556] & 0.537 & 0.000 & 0.996 [0.996, 0.996] & 0.990 [0.988, 0.991] \\
\midrule
Llama-3.2-3B & base & 0.496 & & & & & \\
 & 8 & 0.496 & 0.000 [$-$0.002, 0.002] & 0.003 & 0.003 & 0.022 [0.019, 0.025] & 0.010 [0.008, 0.013] \\
 & 4 & 0.429 & 0.067 [0.057, 0.078] & 0.100 & 0.033 & 0.319 [0.305, 0.332] & 0.307 [0.293, 0.320] \\
 & 3 & 0.054 & 0.442 [0.420, 0.463] & 0.451 & 0.009 & 0.919 [0.910, 0.927] & 0.907 [0.899, 0.915] \\
\bottomrule
\end{tabular}
\end{table}

\paragraph{Results.} Table~\ref{tab:agree} reports the measure. On Qwen2.5-1.5B at 4 bits,
accuracy falls 9.5 points while 40.9\% of predictions change. On Mistral-7B a 4.3-point
accuracy drop comes with a change in 23.2\% of predictions. At 8 bits on Llama-3.2-3B,
accuracy is unchanged at 0.496, regressions and recoveries are each 0.33\%, and Correctness
Agreement still falls to 0.493, by exactly the regressions. At 8 bits 2.2\% of predictions
change, with an interval of 1.9\% to 2.5\%.

\paragraph{The reference depends on its amplitude.} We repeated $D_0$ over three independent
seed pairs at each of three dither amplitudes. On Llama-3.2-3B it is 1.1\% to 1.3\% at 0.02 of a
step, 1.8\% to 2.0\% at 0.05 and 2.3\% to 2.6\% at 0.1. On Qwen2.5-1.5B the three ranges are 1.1\%
to 1.4\%, 1.8\% to 2.1\% and 2.6\% to 3.0\%, and on Mistral-7B 0.6\% to 0.8\%, 0.9\% to 1.0\% and
1.3\% to 1.4\%. Seeds move $D_0$ by a few tenths of a point and the amplitude moves it by a factor
of two. So $D_0$ is a property of the chosen reference, and the 8-bit disagreement of 2.2\% on
Llama is just below the 2.3\% to 2.6\% that a dither of a tenth of a step produces. We read these numbers
descriptively. Whether near-threshold items explain the 8-bit changes would need their
prediction margins measured, which we have not done.

\section{Inferring the consuming operator}\label{supp:trace}

The results so far name the consumer by hand. Someone knew that keys feed the softmax and that
value and output write the residual stream. The tracer is an attempt to remove that step. Given a model, it is meant to map
every parameter tensor to the operator its output flows into, and from that regime and the
quantization target to a quantization rule. The measurement below shows how far it gets.

\begin{table}[h]
\centering
\caption{The rule depends on both the consuming operator and the quantization target. Only the
entries marked (m) are results of this paper. Entries marked (e) were measured in the earlier
cache study. Entries marked (d) are the tracer's defaults. We have not tested them, they are not
recommendations, and the table should not be read as complete. Below 4 bits the weight rules for the score path
and residual are not stable (Section~\ref{sec:proj:below4}), and at 4 bits they were measured
with the absmax scale and reverse under a clipped scale, on one of two models with
rounding and on both with GPTQ
(Section~\ref{sec:proj:quant}).}
\label{tab:regimes}
\small
\begin{tabular}{p{2.4cm}p{3.0cm}p{4.3cm}p{4.3cm}}
\toprule
regime & consuming operator & rule for weights & rule for cached activations \\
\midrule
score path & $\mathrm{softmax}(QK^\top)$ & lower sensitivity at 4 bits (m, Section~\ref{sec:proj}); for keys at 3 bits, keep the rows with the largest weight error precise (m, two models, one budget) & per channel with a zero-point (e) \\
linear residual & writes the residual stream & higher sensitivity at 4 bits (m, Section~\ref{sec:proj}) & per vector (e) \\
gated selection & top-$k$ router & keep precise unless the error spread is below the margin (m, two routers, one quantizer, Section~\ref{sec:route}) & per channel (d) \\
state decay & linear recurrence & log basis (m, three checkpoints of two architectures, Section~\ref{sec:decay}) & per channel (d) \\
normalization & layer or RMS norm & full precision (d) & full precision (d) \\
unknown & no evidence & per channel with a zero-point (d) & per channel with a zero-point (d) \\
\bottomrule
\end{tabular}
\end{table}

Table~\ref{tab:regimes} lists the regimes and their rules. Four cells of it are measured here,
the weight rules for the score path, the residual, the router and the decay. The rest is the
tracer's default behavior, listed so that a reader can see what it would do. Three front ends assign the regime. A
structural pass reads each module's type and name. A layer-norm module is a normalization whatever
its name. A router gate is selection, and the pass excludes the gate projection of a SwiGLU
feed-forward block, which is not a router. The other two read a graph and no name. One traces the
module symbolically with torch.fx \citep{reed2022torchfx}. The other runs the model once on example
inputs and walks the autograd graph of the output. From each softmax or top-$k$ node it walks back
to the first weight matrix on every path, which is the tensor that produced the scores, and a
top-$k$ outranks a softmax that sits before it. From each scan, and from each product carried
through a loop as $h_t = a\,h_{t-1} + b_t$, it labels the parameters that multiply the state
without crossing a matrix product. Where a graph pass has positive evidence it overrides the
structural guess. A tensor with no evidence keeps the structural label. In a graph-only mode,
which reads no parameter name, it gets a default instead: a normalization is recognized by its
module type, any other linear weight is a residual write, and the rest is unknown and gets the
conservative rule that keeps scale on a coordinate whose consumer is unknown.

\paragraph{Measurement.} The paper's own results do not depend on the tracer. Every consumer in
Sections~\ref{sec:proj}, \ref{sec:route} and~\ref{sec:decay} was named by hand. We measured the
tracer on models whose consuming operators are known by construction, and
Section~\ref{app:tracer} gives the details. Symbolic tracing with torch.fx failed on all eight
library architectures we tried, each time on an input check in the forward pass. The autograd
pass ran on all eight. Used alone, with no parameter name read, it labels all 200 scored tensors
correctly. The truth is read from the architectures' own names, so here the method and the truth
are independent. Only 30 of the 200 labels rest on graph evidence: the 16 query and key matrices
of four decoders, the 8 of the two mixture-of-experts models with their 4 router gates, and
Mamba's 2 decay tensors. The other 170 are defaults, and their being right means that the pass
raised no false alarm on a residual write or a normalization. A ninth architecture, RWKV-4, was
held back while the pass was written, and its truth was fixed before the pass was run on it.
There the graph pass finds both decay tensors and gets 16 of 16. The name pass gets 14, because
it takes a feed-forward layer named \texttt{key} for a score path. On four small modules written
with names that carry no signal, the name pass gets 9 of 13 tensors right and the graph pass 13
of 13.

The tracer is therefore a name pass that fits current architectures, plus a graph pass that needs
one forward pass and no tracer support. Its limits are those of the measurement. The models are
small random-weight builds, since only the module graph matters, and we have not run it on a full
checkpoint. A fused tensor that holds query, key and value, as in GPT-2, gets one label. The
recurrence detector looks for one pattern and was written with Mamba's loop in view, with RWKV-4
as its only blind test. And a label says which operator reads a tensor first, not how sensitive
the tensor is, so a calibration batch is still needed to confirm each regime.

\section{Measurement of the tracer}\label{app:tracer}

The paper's own results do not depend on the tracer. Every consumer in
Sections~\ref{sec:proj}, \ref{sec:route} and~\ref{sec:decay} was named by hand. We measured the tracer on two groups
of models whose consuming operators are known by construction. The first group is nine library
architectures, Llama, Mistral, Qwen2, Gemma, GPT-2, Mixtral, OLMoE, Mamba and RWKV-4, each built
from its configuration at a small size with random weights, since only the module graph and the
parameter names matter. RWKV-4 is a blind case: the autograd pass was written without it, and the
rules that give its truth were fixed before the pass was run on it. The second group is four small
modules written for the test, an attention block, a top-$k$ router, a diagonal recurrence and a
SwiGLU block, each in two namings: the usual names, and names that carry no signal. The graph
passes get one forward pass on random inputs. A tensor is scored only if its consumer is
declared. Biases, embeddings, convolutions and Mamba's step-size projections are not scored, and
GPT-2's fused projection is reported separately. In RWKV-4 the key, value and receptance maps of
the attention block are not scored either: the key enters a normalized exponential weighting over
time, and we do not declare whether that is a score path.

\begin{table}[h]
\centering
\caption{The graph pass reads no parameter name. It matches the name pass on the eight
architectures it was written with and is ahead of it on the blind one. On the library
architectures the truth is read from names, which makes it independent of the graph pass and not
of the name pass. Of the 200 graph-only labels on the eight, 30 rest on graph evidence and 170
are defaults.}
\label{tab:tracer}
\small
\begin{tabular}{lrrrrr}
\toprule
 & & \multicolumn{2}{c}{correct} & \multicolumn{2}{c}{trace succeeds} \\
 & scored & names only & graph only & symbolic & forward pass \\
\midrule
eight library architectures & 200 & 200 & 200 & 0 of 8 & 8 of 8 \\
RWKV-4, blind & 16 & 14 & 16 & 0 of 1 & 1 of 1 \\
four small modules, usual names & 13 & 13 & 13 & 4 of 4 & 4 of 4 \\
four small modules, uninformative names & 13 & 9 & 13 & 4 of 4 & 4 of 4 \\
\bottomrule
\end{tabular}
\end{table}

Table~\ref{tab:tracer} gives the result. Symbolic tracing with torch.fx fails on all nine library
architectures, each time on an input check in the model's forward pass. The autograd pass does
not trace symbolically. It runs the model and reads the graph that the run leaves, so it works on
any model that runs. With no name read, it labels the query and key matrices of every decoder as
score path, including where attention goes through the fused scaled-dot-product operator. It
labels the four Mixtral and OLMoE router gates as selection, although a softmax sits between
each gate and its top-$k$. It labels Mamba's decay tensors as state decay, although Mamba's
recurrence is a Python loop with no scan operator in it. It labels no value, output or
feed-forward matrix as anything but a residual write, and that includes the gate projection of
every SwiGLU block.

On RWKV-4 the pass finds the recurrence and labels the decay parameter of both blocks. The name
pass gets 14 of 16 there, because it labels the feed-forward layer named \texttt{key} as a score
path. Used together, the two passes keep that error: the graph pass has positive evidence only
for the score, selection and decay regimes, so it cannot overrule a wrong name on a residual
write.

The pass also labels three kinds of tensor that have no declared truth. GPT-2's fused projection,
which holds query, key and value in one tensor, is labelled as a score path, which is right for
two thirds of it. The bias of Mamba's step-size projection and RWKV-4's \texttt{time\_first}
parameter are labelled as state decay, since each enters the carried product elementwise. We do
not score them.

The second group tests the pass without names. With uninformative names the name pass gets 9 of
13 tensors right: it cannot see the query, key, gate or decay tensors and labels them by its
default. The graph pass gets 13 of 13 in both namings. An earlier version of the tracer had two
defects that this one removes. Its walk back from a scan labelled every linear layer on the way,
so the input map of the recurrence was taken for a decay. The walk now labels only the parameters
that multiply the state. And a name rule written for a fused attention projection labelled any
parameter named \texttt{in\_proj} as a score path. The tracer now also reports, for each graph
pass, whether it failed, ran and found no sink, or ran and found one. The record was made with
transformers 4.46.3. With version 5.5.0, which names the expert weights of the two
mixture-of-experts models differently, 128 tensors are scored on the eight architectures and the
graph pass again gets all of them and 16 of 16 on RWKV-4.

\section{The pair law under synthetic noise}\label{app:law}

Section~\ref{sec:pullback} gives the flip rate under independent
Gaussian noise, counting only the pair at the boundary, as the average of
$\Phi(-m/(\sqrt2\,\sigma))$ over the margins. For small $\sigma$ this is the linear law
of~\eqref{eq:tube}. The rerun stores the margin of every decision, so the average is computed
exactly and no estimate of the density is needed. At $\sigma = 0.01\,m_{50}$, which is 0.06 to
0.07 of the 10th percentile of the margin, the measured flip rates at the top-$k$ boundary are
0.391\% on Mixtral and 0.423\% on OLMoE, against 0.390\% and 0.419\% from the pair law. At the
top-1 boundary of the same two routers they are 0.325\% and 0.386\% against 0.334\% and 0.384\%.
Doubling the noise to $0.02\,m_{50}$ multiplies the measured rate by 1.95 and 1.97 at the
top-$k$ boundary and by 2.06 and 2.01 at the top-1 boundary. Up to $\sigma = 0.1\,m_{50}$ the
measured rate stays within 7\% of the pair law in all four cases. A level at
$0.01\,m_{50}$ rests on about 2{,}000 flips for Mixtral and 4{,}400 for OLMoE, so differences of
2\% to 3\% there are within sampling error.

Above that the measured rate pulls away from the pair law, since pairs other than the one at the
boundary come within reach. At $\sigma = m_{50}$ it is 1.44 times the pair law on Mixtral and
1.49 times on OLMoE, and at $2\,m_{50}$ it is 1.74 and 1.78 times. So the law holds where it
should, for noise well below the typical margin, and fails in the direction and at the scale the
single-wall approximation predicts. This is a test under synthetic independent noise. A
quantization error is not independent across experts, and for it the measured quantity is $R$.

\section{State decay: the clamp of the 32-bit grid, and grids on the time constant}\label{supp:decay}

This section supplements Section~\ref{sec:decay}. The linear grid on $A$ described there has 8
uniformly spaced levels from a tensor's smallest to its largest value of $A$. Two implementations of the same
grid, on the same windows, give $1.0\times10^{10}$ and $5.9\times10^{9}$. They differ in whether
the grid's endpoints and step are held in 64-bit or 32-bit floating point. The top level is
computed as a large negative offset, of order $10^{4}$ in a typical layer and $10^{6}$ in the
extreme one, plus a multiple of the step, which cancels to a value
near zero, so its sign is decided by rounding. It lands at or above zero, where the harness clamps
the decay rate to a small positive floor, in 6 of the 48 layers in one implementation and in 8
in the other. The perplexity of the collapsed model depends on those bits, and neither figure is
a stable number. The count of 8 is in the record of the endpoint-exact run of Section~\ref{sec:decay}. The count of
6 comes from a recomputation and is not in a committed record. The slowdown figures of the 32-bit grid share the dependence: in exact arithmetic the top
level would be the tensor's slowest rate. 

\paragraph{A linear grid on the time constant.} Both linear grids above are uniform in the rate,
so their range is set by the fast channels and the slow ones fall below one step. The opposite
linear grid is uniform in the time constant $\tau = 1/|A| = \exp(-A_{\log})$. Its range is set by
the slow channels, and the fast ones share its lowest level. With one range per tensor, this grid
puts 92\% of the channels of Mamba-790m and 95\% of those of Mamba-370m on its lowest level and
speeds up nearly all of them by more than a factor of 10. No channel slows by a factor of 2. The
perplexities are 71.3 and 162.5, which are 6.1 and 11.6 times the baselines, with intervals of
5.9 to 6.4 and 10.9 to 12.4. That is a large
loss and not a collapse. The two per-tensor linear grids make errors of about the same size in
opposite directions. One slows more than 90\% of the channels tenfold and the perplexity passes
$10^{9}$. The other speeds up more than 90\% of them tenfold and the perplexity rises by an order
of magnitude. Slowing a channel moves $a$ toward 1, where the gain diverges, and speeding it up
moves $a$ away. The asymmetry is what the gain predicts, though we have not isolated it from the
other differences between the two grids.

With one range per row, the grid on $\tau$ gives 13.60 on Mamba-790m and 16.48 on Mamba-370m,
rises of 17\% and 18\%, with intervals of 16\% to 18\% and 17\% to 19\%. It puts 46\% and 47\% of the channels on its lowest level and speeds up
9\% and 13\% of them by more than a factor of 10. Among the slowest tenth of the channels, 0.4\%
and 0.6\% change by more than a factor of 2. So these models tolerate large errors in their fast
channels when the slow ones are kept. The rises are smaller than the per-tensor log grid's 24\%
and 35\%, but the comparison is not at equal cost. A row has 16 entries, and a range per row
stores two endpoints for every 16 three-bit values. A log grid with a range per row, at that same cost, gives 11.91 and 14.27. Those are
rises of 2.3\% and 2.0\%, with intervals of 2.1\% to 2.5\% and 1.8\% to 2.2\%, and the grid
changes 0.04\% and 0.13\% of the channels by more than a factor of 2. So at equal cost the log
basis is the better one by a wide margin, 2\% against 17\% and 18\%. The collapse reported above is a
property of grids uniform in the rate, not of linear grids in general.

\section{Formal checks, and notes to the tables of the paper}\label{supp:checks}

A Lean 4 file, checked against Mathlib with no proof left open, verifies some of the identities
the text relies on: the relations among the agreement counts, the cancellation of a row constant by
the softmax, the split of a score error into row-constant and centered parts, the sufficient
conditions for a top-$k$ set to survive, the invariance of routing to a common scale, the two
facts about subspaces to which Whitney's conditions reduce for the arrangement of ties, an
inclusion for (a) and closure under limits for (b), and the
identities of the decay parameterization. It does not check the integral in
\eqref{eq:tube}, the equivariance of the min-max grid behind the scale invariance of $R$, or the
derivative of the gain, for which it proves the finite-difference identity. It makes no statement about any trained model.

The captions of four tables of the paper are shortened there. Their full text follows.

\emph{Table~\ref{tab:proj}, the matched-bit ratios.} At matched bits the relative Frobenius errors of the four projections are similar, while value and output move the output distribution more. A ratio below 1 means query and key move the output less than value and output. Each ratio is the mean over Q and K divided by the mean over V and O. Gemma was not run at 2 bits. Brackets are 95\% percentile intervals from a paired bootstrap over whole sequences, 2{,}000 resamples. The ratios and intervals come from a second run on the same sequences that recorded each sequence's KL. Every ratio of the first run lies inside its interval. The Frobenius ratios are from the first run.

\emph{Table~\ref{tab:factorial}, the factorial.} The weight-statistic ordering follows the bit allocation. Each entry is the mean over query and key divided by the mean over value and output, so a value above 1 says query and key are the more distorted or the more damaging. With the shipped allocation these two weight statistics put query and key first under both schemes. With equal bits the Frobenius ratio is 1.00 to 1.09 and the kurtosis ratio 1.0 to 1.9. The functional ratio is far below 1 under K-quant at equal type. Runs on 32 sequences. Brackets are 95\% intervals from a paired bootstrap over sequences, from a rerun on Llama-3.2-3B that reproduces every value of the first run. Mistral-7B has no intervals.

\emph{Table~\ref{tab:wave}, the wavelength probe.} Much of the 3-bit key damage on Llama-3.2-3B is removed by sparing the long-wavelength rotary rows. Brackets are 95\% intervals from a paired bootstrap over 32 sequences, from a rerun on Llama-3.2-3B that reproduces the first run. The control has no intervals. Each rank correlation is over eight octiles, and the control's 0.571 is not distinguishable from zero. The coupling of damage to wavelength appears on both models, and its size differs by a factor of 25. Shares are of the summed single-octile KL, which on Llama is far below the joint KL because the damage is superadditive. The spared share is of the joint KL. Row weight error is nearly uncorrelated with wavelength, which is a statement about the two orderings of rows and not about which rows to protect (Table~\ref{tab:protect}).

\emph{Table~\ref{tab:route}, the routing study.} Flips concentrate on low-margin tokens on both routers. Margins and noise are in logits. The figure in parentheses is the flip rate of tokens below the median margin divided by that of tokens above it. At the two smallest noise levels no token above the median flips, so the ratio is not defined. Under 4-bit gate quantization the concentration survives on Mixtral, while on OLMoE nearly every token's expert set changes. The ratio $R$ of error spread to margin separates the two. In the noise rows the second figure is the pair law, the average of $\Phi(-m/(\sqrt2\,\sigma))$ over the measured margins. Single runs.

\bibliographystyle{IEEEtranN}
\bibliography{references}

\end{document}
